\clearpage\thispagestyle{empty}
\pdfbookmark[0]{Quellen zur Besprechung von Humboldt und Oltmans}{humboldt1808-sources}
\begin{center}\textbf{Quellen zur Besprechung von Humboldt und Oltmans}\end{center}
\noindent Diese KI-Übertragung ins heutige Deutsch folgt Gauß’ Besprechung auf S.~509 der \href{https://gdz.sub.uni-goettingen.de/id/PPN236007467}{\textit{Werke}, Bd.~VI} (Göttingen, 1874). Sie erschien ohne Namensnennung unter der Überschrift \emph{Paris und Tübingen} in den \textit{Göttingischen gelehrten Anzeigen}, Stück~193, 3. Dezember 1808, \href{https://api.digitale-sammlungen.de/iiif/image/v2/bsb10538674_00922/full/1800,/0/default.jpg}{S.~1928}. Das digitalisierte Exemplar stammt aus der Bayerischen Staatsbibliothek.
\par\smallskip\noindent Der lateinische Titel bezeichnet eine \emph{Übersicht der geografischen Längen und Breiten, die Alexander von Humboldt in den Jahren 1799 bis 1804 in der Äquatorialregion astronomisch beobachtete und Jabbo Oltmans berechnete}. Gauß nennt F. Schöll und J. G. Cotta in Paris und Tübingen und gibt den Umfang mit sechzehn Folioseiten an. KI-Übertragung und KI-Satz.
\clearpage
